\chapter{Appendix: Scenario Eta - Coppice}

\section{The Legacy Paradigm \\\& Its Failures}
In legacy forest and resource management, data was siloed in centralized Geographic Information Systems (GIS) operated by distant bureaucracies. Field workers relied on ruggedized tablets that became useless off-grid, requiring data to be synced at the end of the day. The interfaces were disconnected from the ecological reality, presenting nature as a spreadsheet of harvestable assets rather than a complex, living system.

\section{The Incremental Automation Pathway}
Phase 1 involved manual logging of tree health and harvest yields into disconnected devices. Phase 2 introduces AR spatial overlays that project ecological data directly onto the forest canopy, assisting stewards in identifying trees ready for coppicing. Phase 3 relies on ambient sensor meshes: the forest itself acts as a living database, communicating soil health and growth rates directly to the local mesh, requiring human intervention only for the physical harvest.

\section{The Human Boundary (Limits of Automation)}
While growth tracking and yield prediction are automated, the decision to actually cut a tree remains a manual, human choice. The interface will highlight optimal trees for coppicing based on ecological algorithms, but it demands physical confirmation from the steward. The system cannot authorize automated harvesting machinery; the human must swing the axe or operate the saw, ensuring physical stewardship and connection to the ecosystem.

\section{Interface Mechanics (The "How")}
Ecological data is stored locally on field nodes and synced opportunistically using CRDTs when stewards cross paths in the forest. AR overlays provide spatial UI without requiring physical screens in the wild. Because mobile nodes have strict battery constraints in the field, thermodynamic rendering ensures that AR graphics seamlessly degrade to minimal haptic feedback or simple audio cues as power reserves dwindle.
